% Figura: soluciones de la ecuación de Airy (datos: scripts/generar_datos.py)
\begin{figure}[ht]
\centering
\begin{tikzpicture}
  \begin{groupplot}[group style={group size=1 by 2, vertical sep=1.5cm},
      oscuro, width=0.92\textwidth, height=5.3cm,
      xmin=-9, xmax=2.4, xtick={-8,-6,-4,-2,0,2}, unbounded coords=jump]
    % ---------- (a) sistema fundamental ----------
    \nextgroupplot[ymin=-1.6, ymax=3.2, ytick={-1,0,1,2,3},
        title={(a) Sistema fundamental $y_1$, $y_2$ (series \eqref{eq:airy-y1y2})},
        legend pos=north west]
      \fill[grisoscuro, opacity=0.45] (axis cs:0,-1.6) rectangle (axis cs:2.4,3.2);
      \addplot[c1, restrict y to domain=-10:10] table[x=x, y=y1] {datos/airy.dat};
      \addplot[c2, restrict y to domain=-10:10] table[x=x, y=y2] {datos/airy.dat};
      \legend{$y_1$, $y_2$}
      \node[font=\scriptsize, text=gris, anchor=south] at (axis cs:-6.6,-1.58) {región oscilatoria ($x<0$)};
      \node[font=\scriptsize, text=gris, anchor=south] at (axis cs:1.2,-1.55) {crecimiento};
    % ---------- (b) funciones de Airy ----------
    \nextgroupplot[ymin=-0.65, ymax=1.6, ytick={-0.5,0,0.5,1,1.5},
        xlabel={$x$},
        title={(b) Funciones de Airy $\Ai$ y $\Bi$, combinaciones \eqref{eq:Ai-Bi}},
        legend pos=north west]
      \fill[grisoscuro, opacity=0.45] (axis cs:0,-0.65) rectangle (axis cs:2.4,1.6);
      \addplot[c3, restrict y to domain=-10:10] table[x=x,
          y expr={0.355028054*\thisrow{y1}-0.258819404*\thisrow{y2}}] {datos/airy.dat};
      \addplot[c4, restrict y to domain=-10:10] table[x=x,
          y expr={1.732050808*(0.355028054*\thisrow{y1}+0.258819404*\thisrow{y2})}] {datos/airy.dat};
      \legend{$\Ai$, $\Bi$}
  \end{groupplot}
\end{tikzpicture}
\caption{Ecuación de Airy $y''=x\,y$. Para $x<0$ las soluciones oscilan con
frecuencia creciente y amplitud decreciente; para $x>0$ (zona sombreada)
crecen exponencialmente, salvo $\Ai$, que decae. Las curvas se han obtenido
sumando directamente las series \eqref{eq:airy-y1y2}, que convergen en todo
$\R$.}
\label{fig:airy}
\end{figure}
